\section{Part II: Linear Algebra with Tensor Networks}
\label{sec:part2}
\lecture{29}{09}{2026}

% Mirrors slides.qmd "Part II" outline. One subsection per outline item.

Part~\partref{1} carried out every operation with two primitives: apply an
operator, then round. That is enough for an explicit scheme and nothing else.
The moment a method requires \emph{solving} rather than \emph{applying}, we need
genuine linear algebra on the MPS manifold.

Part~\partref{2} supplies the two missing capabilities and spends them
immediately. First, solving $Ax = b$ in tensor format
(\cref{subsec:linear-solvers}), which buys unconditional stability for the heat
equation (\cref{subsec:heat-implicit-euler}) and returns in Part~\partref{3} as
the pressure-Poisson step. Second, constructing an MPS from a function we can
only sample (\cref{subsec:finding-qtt}), which is what lets us put an aircraft
wing on a $2^{30}$-cell mesh without touching $2^{30}$ cells
(\cref{subsec:tt-cross-geometry}). \Cref{subsec:part2-open-problems} records
what remains unsolved.

\subsection{Linear System Solvers}
\label{subsec:linear-solvers}
% Keywords: solving $Ax=b$ in TT/QTT format, ALS (alternating least
% squares), DMRG-style two-site sweeps, AMEn (alternating minimal energy),
% conjugate gradient / GMRES in tensor format, convergence criteria,
% preconditioning.

The task: $A$ is given as an MPO, $b$ as an MPS, and we want $x$ as an MPS of
bond dimension $\chi$.

\begin{itemize}
  \item Direct inversion is meaningless here. $A^{-1}$ has no reason to be
    low-rank even when $A$ and $x$ both are. (The QTT Laplacian is the
    instructive exception: its inverse \emph{does} admit an explicit low-rank
    representation~\cite{kazeev2012laplace}, which is precisely why that case is
    remarked upon.)
  \item Instead \textbf{recast the problem as a minimization over the MPS
    manifold},
    \begin{equation}
      \min_x \norm{Ax - b}^2
      \qquad\text{or}\qquad
      \min_x \tfrac{1}{2}\, x^{\top}\! A x - x^{\top} b ,
      \label{eq:tt-linear-min}
    \end{equation}%
    the second form being available when $A \succ 0$, i.e.\ symmetric positive
    definite. Prefer it when you can: it avoids squaring the condition number,
    which the residual form does.
  \item The manifold of fixed-rank tensor trains is not a vector space, so we
    cannot optimize all tensors at once. We optimize \textbf{one tensor at a
    time}, which is a linear problem in that tensor alone.
\end{itemize}

\paragraph{Sweeping solvers: ALS and DMRG.}
\margintext{The one-site/two-site distinction is the same one that separates
White's original DMRG from its later single-site variants. The trade is rank
adaptivity against local cost.}
\begin{itemize}
  \item \textbf{ALS, or one-site DMRG.} Freeze every tensor but one. With the
    chain in mixed canonical form centred on that tensor, the problem collapses
    to a small dense linear system of size $\chi^2 d$, solved exactly. Sweep
    left to right and back~\cite{holtz2012als}.
  \item Each local solve is exact, so the global objective decreases
    monotonically. Convergence to a global optimum is not guaranteed; ALS can
    stall in a local minimum of the fixed-rank manifold.
  \item Its limitation: $\chi$ is fixed in advance and never adapts.
  \item \textbf{Two-site DMRG.} Optimize a \emph{pair} of neighbouring tensors,
    then split them by SVD. The truncation tolerance in that SVD sets the new
    bond dimension, so $\chi$ adapts on its own. The local problem is now of
    size $\chi^2 d^2$ and the local cost rises to $O(\chi^6)$.
  \item Two-site updates also escape local minima that trap one-site
    sweeps, because the pair carries variational freedom that neither tensor has
    alone.
\end{itemize}

\paragraph{AMEn and Krylov alternatives.}
\begin{itemize}
  \item \textbf{AMEn}~\cite{dolgov2014amen} --- one-site sweeps enriched with
    directions taken from the current residual $b - Ax$. This recovers rank
    adaptivity at roughly one-site cost, and is often the best default.
  \item \textbf{Krylov methods in TT format} --- CG or GMRES with a rounding
    step after every operation. Easy to build on top of existing TT arithmetic.
  \item But rank growth per iteration is severe: each matrix--vector product
    multiplies $\chi$ by $\chi_{\mathrm{MPO}}$ and each vector addition adds
    ranks, so truncation must be aggressive, and aggressive truncation destroys
    the Krylov subspace's orthogonality.
  \item Preconditioning in TT form remains open
    (\cref{subsec:part2-open-problems}).
\end{itemize}

\begin{remark}[Rule of thumb]
  Use sweeping methods for structured PDE operators, where the MPO is known in
  closed form and warm starts are available. Use Krylov methods when the
  operator is effectively a black box.
\end{remark}

\paragraph{What the solve actually costs.}
This is not a peripheral concern. In our incompressible CFD solver the
DMRG-type Poisson solve dominates every other part of the
pipeline~\cite{peddinti2024cfd}:

\begin{center}
  \begin{tabular}{ll}
    \toprule
    \textsc{Quantity} & \textsc{Value} \\
    \midrule
    Share of total runtime  & $80.1\%$ \\
    Worst-case local cost   & $O(n\chi^6)$ \\
    Observed total cost     & $O(n\chi^{4.1})$ \\
    \bottomrule
  \end{tabular}
\end{center}

\noindent
\margintext{The measured exponent $4.1$ against a worst case of $6$ is where all
the practical engineering lives. Warm-starting from the previous timestep is the
single largest contributor.}
The gap between the worst case and what is measured is where the practical knobs
live: sweep count, local tolerance, rank cap, and above all
\textbf{warm-starting from the previous timestep}. The worst-case $\chi^6$
assumes the local eigen- or linear solve runs to full accuracy at maximum rank
on every site of every sweep; warm starts mean it almost never does.

\subsection{Solving the Heat Equation (Implicit Euler Method)}
\label{subsec:heat-implicit-euler}
% Keywords: implicit time-stepping, unconditional stability, linear
% system per time step, QTT operator inversion, comparison with explicit
% scheme, cost per step vs. accuracy/stability trade-off.

Return to the heat equation~\eqref{eq:heat} of
\cref{subsec:heat-explicit-euler}, now with backward Euler:
\begin{equation}
  \left(\mathbb{1} - \alpha \Delta t\, L\right) u^{k+1} = u^{k}.
  \label{eq:heat-implicit}
\end{equation}

\paragraph{Why bother.}
\begin{itemize}
  \item Equation~\eqref{eq:heat-implicit} requires a \textbf{linear solve} every
    step instead of a matrix--vector product: strictly more work per step.
  \item The payoff is unconditional stability. The amplification factor is
    $(1 + \alpha \Delta t\, \lambda)^{-1}$ for every eigenvalue $\lambda \ge 0$
    of $-L$, so it never exceeds one whatever $\Delta t$ is. The step size is
    then set by \textbf{accuracy}, not by $\Delta x^2$.
  \item With $\Delta x = 2^{-n}$ and $n \gtrsim 20$, the explicit bound
    $\Delta t \le \Delta x^2/2\alpha$ is unusable. That is the entire argument,
    and it strengthens with every bit of resolution quantics buys.
\end{itemize}

\begin{remark}
  Note the shape of the trade. Quantics makes very large $n$ affordable in
  memory; the explicit stability condition then makes large $n$ unaffordable in
  time. The compression only pays off in combination with an implicit scheme,
  which is only affordable because of \cref{subsec:linear-solvers}.
\end{remark}

\paragraph{Assembling and solving.}
\begin{itemize}
  \item Build $A = \mathbb{1} - \alpha \Delta t\, L$ directly as an MPO. Its
    rank is $\approx 3$ and the cores are written down
    analytically~\cite{kazeev2012laplace}; there is no numerical construction
    step and no assembly error.
  \item The right-hand side is simply the previous state MPS.
  \item Solve with DMRG or AMEn, \textbf{warm-started from $u^k$}. Consecutive
    solutions differ by $O(\Delta t)$, so one or two sweeps usually suffice.
\end{itemize}

\begin{remark}[How to compare schemes]
  Compare \textbf{total wall-clock time to a target error}, never cost per step.
  Implicit stepping loses on the second measure and wins on the first, because
  it takes far fewer steps.
\end{remark}

\paragraph{The error budget.}
The budget now has three parts rather than two:

\begin{enumerate}
  \item time discretization, $O(\Delta t)$;
  \item space discretization, $O(\Delta x^2)$;
  \item \textbf{the solver residual} $\norm{Au^{k+1} - u^{k}}$ --- new, and
    entirely under your control.
\end{enumerate}

\noindent
Keep the residual well below the two discretization errors. Otherwise you
converge efficiently to the wrong answer, which is the least useful of all
possible outcomes. A practical rule: target a relative residual one order of
magnitude below the expected discretization error, and no tighter, since
sweeping past that point is pure cost.
\margintext{Forward pointer: this exact solve reappears in
\cref{subsec:qtt-cfd-frameworks} as the pressure-Poisson step of Chorin
projection, where it consumes $80\%$ of the runtime.}

\subsection{How to Find the QTT of Your Function?}
\label{subsec:finding-qtt}
% Keywords: constructing QTT from analytic/black-box functions,
% interpolation vs. exact construction, rank growth with resolution,
% sampling on a quantics grid, adaptive rank truncation.

We have been assuming the initial condition arrives already in MPS form. That
assumption has to be discharged, and how it is discharged decides whether the
method scales. There are three routes.

\begin{center}
  \begin{tabular}{@{}lll@{}}
    \toprule
    \textsc{Route} & \textsc{Cost} & \textsc{Requires} \\
    \midrule
    Analytic cores      & free                        & closed form for $f$ \\
    TT-SVD              & $O(2^n)$                    & the dense vector \\
    TT-cross / TCI      & $O(n\chi^2)$ evaluations    & only a black box $f$ \\
    \bottomrule
  \end{tabular}
\end{center}

\begin{itemize}
  \item \textbf{Analytic.} For the families tabulated in
    \cref{subsec:qtt-functions-operators} the small cores are written down by
    hand~\cite{oseledets2013functions}. Exact: no sampling, no error.
  \item \textbf{From a full array (TT-SVD).} Exact and quasi-optimal
    (\cref{subsec:mps}), but it needs the dense $2^n$ vector, so it is viable
    only for small $n$.
  \item \textbf{Black-box (TT-cross, or TCI).} Build the MPS from $O(n\chi^2)$
    \textbf{evaluations} of $f$, never forming the full
    tensor~\cite{oseledets2010ttcross}.
\end{itemize}

The third route is the one that scales, and the one that makes the method
practical rather than merely elegant.
\margintext{TT-cross is what keeps the programme from being circular. Without
it you would need the $2^n$-vector in order to compress the $2^n$-vector.}

\paragraph{Cross interpolation and maxvol.}
\begin{itemize}
  \item Cross interpolation is a \textbf{skeleton (CUR) decomposition applied
    recursively} down the chain.
  \item A rank-$r$ matrix is exactly reconstructible from $r$ well-chosen rows
    and columns. The only question is \emph{which} ones.
  \item \textbf{maxvol}: choose the $r \times r$ submatrix of maximum volume,
    i.e.\ largest $\abs{\det}$. This is near-optimal and computable greedily by
    row swaps.
  \item Quantics combined with cross interpolation gives \textbf{QTCI}:
    parsimonious representations of multivariate functions at very high
    resolution~\cite{ritter2024qtci}. Mature implementations exist in
    \texttt{xfac} and \texttt{Tensor\-Cross\-Inter\-polation.jl}~\cite{nunez2025tci}.
\end{itemize}

The algorithm is spelled out in \cref{subsec:appendix-tt-cross}.

\paragraph{TT-cross in practice.}
The working recipe is short:

\begin{enumerate}
  \item choose the bit ordering, interleaved or stacked
    (\cref{subsec:qtt-functions-operators});
  \item set a tolerance;
  \item sweep until the cross error plateaus;
  \item round.
\end{enumerate}

\begin{remark}[Failure modes to expect]
  \begin{itemize}
    \item A discontinuous $f$ causes rank blow-up: a step edge has singular
      values decaying like $1/k$, so no finite $\chi$ captures it well.
    \item Poorly placed pivots cause features to be missed entirely. The method
      never samples them, so it never learns they exist, and the reported cross
      error stays small while the answer is wrong.
    \item A noisy $f$ is interpolated faithfully, noise included. Cross
      interpolation is exact on the fibres it samples, which is a virtue for
      clean data and a liability for dirty data.
  \end{itemize}
\end{remark}

\subsection{Encoding Geometry via TT-Cross}
\label{subsec:tt-cross-geometry}
% Keywords: TT-cross / tensor cross interpolation (TCI), maxvol algorithm,
% black-box tensor construction without full evaluation, representing
% complex domains/boundaries in QTT, indicator functions, adaptive
% sampling cost.

Now the payoff, and the first place where the tensor-network constraints force a
genuine modelling decision rather than an implementation detail.

\paragraph{The problem with sharp boundaries.}
\begin{itemize}
  \item The goal is to place an arbitrary immersed object --- a cylinder, an
    airfoil --- on a $2^n$ mesh \textbf{without touching $2^n$ cells}.
  \item The natural object is an indicator function $\theta(q(x))$: one inside
    the body, zero outside, where $q$ is a level-set function vanishing on the
    boundary.
  \item That is exactly the wrong thing for a tensor network. A step edge has
    \textbf{slowly decaying singular values}, so $\chi$ explodes. This is the
    failure mode of \cref{subsec:mps} in its purest form: the indicator is
    maximally non-smooth exactly where the physics happens.
\end{itemize}

The fix is to give up the sharp edge \emph{deliberately}, and then bound what
that costs.

\paragraph{Smoothed masks.}
\margintext{This construction is from our 2024 framework. The point of the
$\sqrt{\ell_0^{D-2}/\alpha}$ bound is that the approximation is a modelling
choice with an error bar, not an undocumented hack.}
Replace the indicator by a smooth mask, with $q$ vanishing on the object
boundary $\partial\Omega$~\cite{peddinti2024cfd}:
\begin{equation}
  m(x) = 1 - e^{-\alpha\left(q(x) + \abs{q(x)}\right)},
  \qquad q|_{\partial\Omega} = 0 .
  \label{eq:smoothed-mask}
\end{equation}

\begin{itemize}
  \item The combination $q + \abs{q}$ vanishes identically wherever $q < 0$, so
    $m$ is exactly zero inside the body and rises smoothly outside it. No
    approximation is made in the interior.
  \item $\alpha$ trades sharpness against rank: large $\alpha$ approaches the
    true indicator and costs bond dimension.
  \item The $\ell_2$ distance to the true indicator scales as
    $\sqrt{\ell_0^{D-2}/\alpha}$ in $D$ dimensions, where $\ell_0$ is the
    boundary length scale. The error is therefore controlled by a single knob
    and decays as $\alpha^{-1/2}$.
\end{itemize}

\paragraph{Putting it together.}
\begin{itemize}
  \item Build $m$ as an MPS by TT-cross to tolerance $\varepsilon$, at a cost
    independent of the mesh size.
  \item Apply it by an \textbf{elementwise (Hadamard) MPS product} on the
    velocity field each step. No-slip is enforced without ever meshing the body,
    and without an immersed-boundary force term.
  \item Inlet and outlet conditions fold into the differential MPOs via
    \textbf{ghost cells}, so open boundary conditions keep the domain small.
  \item Demonstrated on a \textbf{NACA 0040 airfoil}: $\chi = 45$, $18{,}800$
    parameters, a compression factor of $27.9$.
\end{itemize}

\subsection{Open Problems}
\label{subsec:part2-open-problems}
% Keywords: rank blow-up for non-smooth data, robustness of TT-cross,
% error control, nonlinear operators in TT format, scalability to 3D/time.

Worth being explicit about what is not settled. The honest list is short enough
to state and long enough to matter.

\paragraph{Ranks and nonlinearity.}
\begin{itemize}
  \item \textbf{Rank blow-up.} Shocks, sharp interfaces, non-smooth data and
    high-$\mathrm{Re}$ turbulence all push $\chi$ up. Where exactly does the
    method stop paying for itself?
  \item \textbf{Nonlinearity.} Products of MPSs square the bond dimension.
    Picard iteration converges linearly and Newton quadratically, but every
    iteration re-rounds, and the interaction between iteration error and
    truncation error is not understood.
  \item \textbf{Error control.} There is no sharp a-posteriori bound linking the
    per-step truncation tolerance $\varepsilon$ to the PDE error accumulated
    over many steps. In practice one assumes linear accumulation, $O(K\varepsilon)$
    after $K$ steps, and checks it empirically.
  \item \textbf{Robustness of TT-cross.} Pivot selection is heuristic, with no
    guarantee for adversarial or noisy black boxes.
\end{itemize}

\paragraph{Solvers and structure.}
\margintext{Of everything on this page, preconditioning has the largest
immediate payoff. Recall that the Poisson solve is $80\%$ of runtime.}
\begin{itemize}
  \item \textbf{Preconditioning in TT format} is largely unsolved, and it is the
    single biggest lever on solver cost. The difficulty is structural: a good
    preconditioner approximates $A^{-1}$, and low-rank $A$ rarely has low-rank
    inverse.
  \item \textbf{Conservation.} SVD truncation preserves neither mass nor energy
    nor $\nabla\cdot\mathbf{v} = 0$; it is optimal in the Frobenius norm, which
    knows nothing about the physics. Structure-preserving rounding is wide open.
  \item \textbf{Scaling out.} Three dimensions plus time, index ordering at
    scale, memory-bound implementations, and distributed or GPU sweeps.
  \item \textbf{Fair benchmarking.} Honest comparison against tuned classical
    CFD --- multigrid, spectral methods --- on equal hardware.
\end{itemize}
